\subsection{Execution Algorithms}

\begin{principle} \hlt{Algo as Layered Decision System}\\
An execution algorithm is not only a schedule. It is a layered system that turns a parent order into child orders, adapts to market state, and records enough context to explain the realised cost afterwards.
\end{principle}

\begin{taxonomy} \hlt{Algorithm Layer Segmentation}
\begin{enumerate}[label=\roman*.]
\setlength{\itemsep}{0pt}
\item \hlt{execution parameters}: side, size, benchmark, timing, limits, participation, urgency, and venues;
\item \hlt{risk-control layer}: sets the curve from impact, timing risk, alpha, and context;
\item \hlt{tactical layer}: uses SORs, liquidity seeking, pegging, and dark probes to realise it.
\end{enumerate}
Risk control sets the envelope; tactics work inside it.
\end{taxonomy}

\begin{principle} \hlt{Benchmark Customisation}\\
Execution styles can be combined: a VWAP can carry minimum participation, limit-price, or would-level constraints; a PoV can be capped by urgency; an IS schedule can still use an SOR and dark-pool probes. The benchmark defines the reference, while constraints define acceptable behaviour around it.
\end{principle}

\begin{caveat} \hlt{Hard Mixtures Are Fragile}\\
Binary benchmark switches are brittle in noisy markets. A trading envelope gives the schedule room to adapt without abandoning risk control.
\end{caveat}

\begin{definition} \hlt{Trading Envelope}\\
A trading envelope is a band of acceptable execution trajectories around the target schedule. Being too fast implies pausing or posting passively; being too slow implies crossing the spread or increasing participation.
\end{definition}

\begin{taxonomy} \hlt{Envelope Construction Segmentation}
\begin{enumerate}[label=\roman*.]
\setlength{\itemsep}{0pt}
\item \hlt{benchmark path}: VWAP, TWAP, PoV, IS, or minimum-impact curve;
\item \hlt{risk width}: tolerance for timing risk, price drift, and completion uncertainty;
\item \hlt{market inputs}: volume curve, volatility, spread, depth, and impact estimates;
\item \hlt{constraints}: limits, participation caps, venue restrictions, and common portfolio end times;
\item \hlt{tactics}: liquidity capture, dark probes, pegging, and routing choices used inside the band.
\end{enumerate}
\end{taxonomy}

\begin{process} \hlt{Customisation Workflow}
\begin{enumerate}[label=\roman*.]
\setlength{\itemsep}{0pt}
\item choose the desired risk profile and inject hard constraints into the envelope;
\item choose liquidity tactics consistent with investment style and observed market state;
\item monitor whether the realised path remains inside the envelope.
\end{enumerate}
\end{process}

\begin{process} \hlt{Execution Decision Loop}
\begin{enumerate}[label=\roman*.]
\setlength{\itemsep}{0pt}
\item initialise the parent order and benchmark;
\item compute a target trajectory from historical profiles and model parameters;
\item observe quotes, trades, fills, rejects, cancellations, and venue responses;
\item revise child-order size, price, urgency, and venue split;
\item publish monitoring and performance indicators for the trader.
\end{enumerate}
\end{process}

\begin{taxonomy} \hlt{Pre-Trade Segmentation}
\begin{enumerate}[label=\roman*.]
\setlength{\itemsep}{0pt}
\item \hlt{expected cost}: spread, impact, volatility, and typical quantity on the book;
\item \hlt{profile}: expected volume curve, duration, and participation bounds;
\item \hlt{context}: news, volatility regime, sector, country, currency, and event risk;
\item \hlt{portfolio coupling}: common end times, exposure constraints, and cross-order synchronisation.
\end{enumerate}
\end{taxonomy}

\begin{taxonomy} \hlt{Order-Submission Control Segmentation}
\begin{enumerate}[label=\roman*.]
\setlength{\itemsep}{0pt}
\item \hlt{urgency}: cross, post, peg, or wait;
\item \hlt{price}: limit level, discretion range, and protection price;
\item \hlt{quantity}: child size, minimum fill, reserve display, and parent slicing;
\item \hlt{exposure}: lit, dark, brokered, indicated, or anonymous;
\item \hlt{intermediation}: single broker for control or multiple brokers for concealment;
\item \hlt{venue}: one book, smart-routed books, crossing networks, or upstairs negotiation.
\end{enumerate}
The execution algorithm controls transaction cost through these levers, not only through the target schedule.
\end{taxonomy}

\begin{principle} \hlt{Reactive vs Proactive Execution}\\
Reactive execution waits for liquidity to appear and usually lowers leakage. Proactive execution advertises interest or crosses the spread to create a trade. The right point on this continuum depends on urgency, search cost, and how damaging exposure would be.
\end{principle}

\begin{taxonomy} \hlt{Portfolio Trading Segmentation}
\begin{enumerate}[label=\roman*.]
\setlength{\itemsep}{0pt}
\item \hlt{program trade}: simultaneous orders across many names, often index or factor related;
\item \hlt{basket trade}: trade a list as one risk package;
\item \hlt{package trade}: dealer quotes a firm bid or offer for the whole portfolio;
\item \hlt{portfolio transition}: move from one target portfolio to another under tracking-error and cost limits.
\end{enumerate}
Portfolio trading reduces leg count and search cost, but transfers basket characterisation and residual-risk pricing to the dealer or algorithm.
\end{taxonomy}

\begin{mechanism} \hlt{Portfolio Insurance as Dynamic Execution}\\
Portfolio insurance synthetically buys downside protection by dynamically reducing risky exposure as the portfolio falls and increasing exposure as it rises. It transfers risk only if the market can absorb the hedge trades; when many investors run the same rule, the execution demand itself becomes market impact.
\end{mechanism}

\begin{proposition} \hlt{Dynamic Hedging Feedback}\\
Let a rule target risky holdings $q(S)$ as a function of market level $S$, with $q'(S)>0$. A price fall $dS<0$ forces $dq=q'(S)dS<0$, so the rule sells into the decline. If temporary impact is $\lambda dq$ with $\lambda>0$, the hedge flow has the same sign as the initial price move.
\begin{proof}
Since $q'(S)>0$, $dS<0$ implies $dq<0$. A sell order has negative signed flow and negative temporary price impact, so the induced impact reinforces the initial decline.
\end{proof}
\end{proposition}

\begin{taxonomy} \hlt{Basket Characterisation Inputs}
\begin{enumerate}[label=\roman*.]
\setlength{\itemsep}{0pt}
\item \hlt{risk}: beta, factor loadings, volatility, sector, currency, and country exposures;
\item \hlt{liquidity}: ADV, spreads, depth, index membership, and short-sale or borrow constraints;
\item \hlt{execution}: names, weights, limits, urgency, venue eligibility, and close/open constraints.
\end{enumerate}
For index-like baskets, diversified residual risk lets dealers quote tighter than for unrelated single-stock blocks.
\end{taxonomy}

\begin{taxonomy} \hlt{Index Rebalance Controls}
\begin{enumerate}[label=\roman*.]
\setlength{\itemsep}{0pt}
\item \hlt{timing}: trade before, during, or after the effective window subject to tracking error;
\item \hlt{substitution}: use futures, ETFs, or correlated names when constituents are expensive;
\item \hlt{liquidity supply}: post or cross opportunistically when forced flow pays a premium.
\end{enumerate}
Index discipline constrains the portfolio, but execution still chooses how much immediacy to buy.
\end{taxonomy}

\begin{table}[htbp]
\centering
\scriptsize
\stepcounter{table}
\renewcommand{\arraystretch}{1.14}
\setlength{\tabcolsep}{4pt}
\begin{tabularx}{\textwidth}{@{}>{\RaggedRight\arraybackslash}p{0.14\textwidth} >{\RaggedRight\arraybackslash}p{0.17\textwidth} >{\RaggedRight\arraybackslash}p{0.31\textwidth} >{\RaggedRight\arraybackslash}X@{}}
\hline
\textbf{Benchmark} & \textbf{Stock type} & \textbf{Trade type} & \textbf{Main feature} \\
\hline
PoV & Medium to large market depth & Long-duration position trade & Flow-following; reactive within participation bounds \\
VWAP/TWAP & Any market depth & Hedging order; long-duration position; tracking-error unwind & Volume-curve schedule; passive when possible \\
IS (Implementation Shortfall) & Medium liquidity depth & Alpha extraction; non-linear hedge, typically gamma; inventory-driven trade & Fast on opportunity; cuts losses on adverse moves \\
Liquidity seeker & Poor fragmented market depth & Alpha extraction; opportunistic position mounting; already split or scheduled order & Pool search; dispersed liquidity; low signalling \\
\hline
\end{tabularx}

\vspace{0.25em}
{\footnotesize Table \thetable. Typical features of trading algorithms.}
\end{table}

\begin{practitionernote} \hlt{What-If Controls}\\
Useful tools support live counterfactuals: finish now, slow down, switch benchmark, seek dark liquidity, or cap participation. The aim is disciplined comparison before changing trajectory.
\end{practitionernote}
